%======================================================================
\section{Interface sem fio}\label{sec:semfio}
%======================================================================
O sistema tem dois enlaces sem fio, os dois em 2,4\,GHz no rádio do próprio ESP32-S3: o
\textbf{ESP-NOW}, da pulseira ao hub, e o \textbf{Wi-Fi}, do hub ao ponto de acesso da
ILPI, por onde passa o MQTT até a central.

\subsection{Modos do ESP32-S3 e o papel de cada nó}\label{sec:modos}

O rádio do ESP32-S3 pode operar como estação (STA, associada a um ponto de acesso), como
ponto de acesso (SoftAP), nos dois ao mesmo tempo (AP+STA) e com o ESP-NOW, um protocolo da
Espressif que troca quadros curtos direto entre dispositivos, na camada de enlace, sem
associação~\cite{espnow}. O ESP-NOW funciona sobre a interface STA ou SoftAP, no canal em que
o rádio estiver.

\begin{table}[H]
\centering
\caption{Modos do rádio em cada nó.}
\label{tab:modos}
\small
\begin{tabular}{@{}l>{\raggedright\arraybackslash}p{4.0cm}>{\raggedright\arraybackslash}p{7.8cm}@{}}
\toprule
Nó & Modo & Por quê \\
\midrule
Hub & STA associado ao Wi-Fi da ILPI, sem economia de energia, + ESP-NOW & precisa do Wi-Fi para o MQTT; sem economia de energia, porque com ela o rádio dorme entre os \textit{beacons} e perde quadros ESP-NOW \\
Pulseira & STA sem associação, só ESP-NOW & o rádio só liga para enviar; não precisa manter a associação nem conhecer a senha da rede \\
Central & --- (rede local, cabo ou Wi-Fi) & --- \\
\bottomrule
\end{tabular}
\end{table}

A primeira proposta, no Projeto 1 deste semestre, já usava o ESP-NOW, mas punha o hub em
AP+STA, com o ESP-NOW num canal fixo, sem criptografia e com mensagens em texto. Ela foi
revista: o ESP-NOW funciona sobre a interface STA, então o SoftAP não acrescentava nada; e o
que evita perder quadros no hub é desligar a economia de energia do rádio. O canal passou a
ser descoberto pela pulseira, e as mensagens ganharam formato binário e criptografia.

\subsection{ESP-NOW: da pulseira ao hub}\label{sec:espnow}

\begin{table}[H]
\centering
\caption{O enlace ESP-NOW.}
\label{tab:espnow}
\small
\begin{tabular}{@{}>{\raggedright\arraybackslash}p{3.0cm}>{\raggedright\arraybackslash}p{10.8cm}@{}}
\toprule
Item & Valor \\
\midrule
Sentido & pulseira envia, hub recebe; sempre \textit{unicast} \\
Endereçamento & o MAC do hub, configurado na pulseira; o hub aceita só a pulseira pareada \\
Criptografia & ligada: chave primária (PMK) comum e chave local (LMK) do par, 16 bytes cada, fora do repositório \\
Canal & o do Wi-Fi do hub, descoberto pela pulseira por varredura \\
Tamanho & até 13 bytes por mensagem (o limite do ESP-NOW é 250) \\
Confirmação & ACK da camada MAC do hub, informado pelo \textit{callback} de envio \\
\bottomrule
\end{tabular}
\end{table}

\paragraph{Formato.} As mensagens são estruturas empacotadas, em \textit{little endian},
definidas num cabeçalho C (\code{protocol\_espnow.h}) que as duas placas compilam
(Figura~\ref{fig:espnow_quadro}). O cabeçalho comum traz a versão do formato, o tipo, um
número de sequência e o tempo desde o boot da pulseira. São seis tipos: \code{PING} (varredura
de canal), \code{HEARTBEAT}, \code{FALL} e \code{HELP} (o pedido de ajuda, só o cabeçalho), e
os cancelamentos \code{FALL\_CANCEL} e \code{HELP\_CANCEL}, que levam a sequência do alerta
cancelado. O hub descarta versão ou tipo desconhecidos e mensagens com tamanho diferente do
esperado para o tipo.

\begin{figure}[H]
\centering
\resizebox{0.95\linewidth}{!}{% Mensagens ESP-NOW: cabeçalho de 8 bytes e o payload de cada tipo (bytes).
\begin{tikzpicture}[font=\scriptsize]
  \tikzset{byte/.style={draw=black!60, fill=#1, minimum height=0.62cm, align=center, inner sep=1pt, anchor=west},
           byte/.default=white}
  \def\u{0.62}
  % cabeçalho comum
  \foreach \y/\nome/\tam in {0/FALL/13, -1.0/HEARTBEAT/12, -2.0/FALL\_CANCEL \textrm{e} HELP\_CANCEL/10, -3.0/PING \textrm{e} HELP/8} {
    \node[anchor=east, font=\scriptsize] at (-0.15,\y) {\code{\nome} (\tam\,B)};
    \node[byte=black!8, minimum width=\u cm] (v) at (0,\y) {ver.};
    \node[byte=black!8, minimum width=\u cm] (t) at (v.east) {tipo};
    \node[byte=corhub!18, minimum width=2*\u cm] (s) at (t.east) {\code{seq}};
    \node[byte=corhub!18, minimum width=4*\u cm] (u) at (s.east) {\code{uptime\_ms}};
  }
  % payloads
  \node[byte=corvest!22, minimum width=4*\u cm] (a2) at ($(0,0)+(8*\u,0)$) {\code{a2\_peak} (float, g$^2$)};
  \node[byte=corvest!22, minimum width=\u cm*1.6] at (a2.east) {\code{prob}\\\%};
  \node[byte=corcentral!22, minimum width=2*\u cm] (bat) at ($(0,-1.0)+(8*\u,0)$) {bateria (mV)};
  \node[byte=corcentral!22, minimum width=2*\u cm] at (bat.east) {atividade (mg)};
  \node[byte=corvest!22, minimum width=2*\u cm] at ($(0,-2.0)+(8*\u,0)$) {\code{seq} do\\[-2pt]alerta};
  % réguas
  \draw[decorate, decoration={brace, amplitude=4pt}] (0,0.42) -- (8*\u,0.42) node[midway, above=5pt] {cabeçalho comum: 8 bytes};
  \draw[decorate, decoration={brace, amplitude=4pt}] (8*\u,0.42) -- (13*\u+0.6*\u,0.42) node[midway, above=5pt] {payload do tipo};
\end{tikzpicture}
}
\caption{Mensagens ESP-NOW da pulseira: cabeçalho de 8 bytes e o conteúdo de cada tipo.}
\label{fig:espnow_quadro}
\end{figure}

\begin{lstlisting}[style=c]
typedef struct __attribute__((packed)) {
    uint8_t  version;     /* ESPNOW_MSG_VERSION */
    uint8_t  type;        /* espnow_msg_type_t */
    uint16_t seq;         /* +1 per new message; a resend repeats it */
    uint32_t uptime_ms;   /* sender uptime */
} espnow_header_t;

typedef struct __attribute__((packed)) {
    espnow_header_t hdr;
    float   a2_peak;      /* filtered |a|^2 at the trigger peak, in g^2 */
    uint8_t prob_pct;     /* the classifier's P(fall), 0-100 % */
} espnow_fall_t;
\end{lstlisting}

\paragraph{Varredura de canal.} A pulseira não está associada a nada e não sabe o canal da
rede. No boot, e depois de 3 falhas seguidas de envio, ela manda um \code{PING} ao hub nos
canais 1 a 13; o primeiro canal que confirma é guardado na memória RTC e usado dali em
diante. Se nenhum responder, ela espera 30\,s e recomeça.

\paragraph{Confiabilidade.} O ESP-NOW não garante a entrega, e uma queda perdida é o pior
erro do sistema. Por isso:
\begin{itemize}
  \item uma mensagem só conta como entregue com o ACK da camada MAC do hub;
  \item \code{FALL}, \code{HELP} e os cancelamentos têm até 10 tentativas, a cada 50\,ms, com o
        mesmo número de sequência. Se todas falharem, a pulseira refaz a varredura e tenta de novo;
        com o MQTT da pulseira ligado, publica o alerta direto no broker;
  \item se nem assim a mensagem chegar, a pulseira guarda o quadro e o reenvia \emph{igual}
        2\,s depois, dobrando a espera até 30\,s, e na hora em que um heartbeat voltar a ser
        entregue. Repete até entregar;
  \item o hub guarda o último par (sequência, \textit{uptime}) de cada pulseira e descarta a
        mensagem que repete os dois: um reenvio cujo ACK se perdeu não vira outro alerta. O
        \textit{uptime} distingue um reenvio da primeira mensagem de uma pulseira que
        reiniciou, e a sequência começa num valor aleatório a cada boot, para o
        identificador do alerta não repetir o de uma queda anterior.
\end{itemize}

\paragraph{Vivacidade.} A pulseira manda um \code{HEARTBEAT} a cada 30\,s. Qualquer mensagem
conta como sinal de vida; depois de 90\,s sem nenhuma, o hub publica a pulseira como
\code{offline} e abre o alerta \code{pulseira\_offline}.

A Figura~\ref{fig:sequencia} junta tudo: o caminho de uma queda até o celular do cuidador,
com as metas de latência, e o cancelamento pelo botão.

\begin{figure}[H]
\centering
\resizebox{\linewidth}{!}{% Sequência de uma queda, da pulseira ao Telegram, e do cancelamento.
\begin{tikzpicture}[font=\scriptsize, >=Stealth]
  \tikzset{ator/.style={bloco=#1, minimum width=2.3cm, minimum height=0.75cm},
           msg/.style={->, thick, draw=#1}, msg/.default=black!75,
           resp/.style={->, thick, dashed, draw=black!45},
           nota/.style={draw=black!30, fill=yellow!8, align=left, font=\tiny, inner sep=2pt, anchor=west}}
  \foreach \x/\n/\t/\c in {0/v/Pulseira/corvest, 3.4/h/Hub/corhub, 6.8/m/Mosquitto/corcentral,
                          10.2/n/notificador/corcentral, 13.6/t/Telegram/cinza} {
    \node[ator=\c] (\n) at (\x,0) {\textbf{\t}};
    \draw[black!30, thick] (\n.south) -- (\x,-8.6);
  }
  % queda
  \node[nota] at (0.1,-0.85) {pico $\ge$ 2\,g $\to$ janela $\to$\\Random Forest: $P \ge 0{,}73$};
  \draw[msg=corvest] (0,-1.6) -- node[rotulo, above] {\code{FALL} (ESP-NOW, cifrado)} (3.4,-1.6);
  \draw[resp] (3.4,-2.0) -- node[rotulo, below] {ACK da camada MAC} (0,-2.0);
  \node[nota] at (3.5,-2.65) {regra \code{queda}:\\LED vermelho};
  \draw[msg=corhub] (3.4,-3.3) -- node[rotulo, above] {\code{alerta} \code{ativo} (QoS 1)} (6.8,-3.3);
  \draw[resp] (6.8,-3.7) -- node[rotulo, below] {PUBACK} (3.4,-3.7);
  \draw[msg=corcentral] (6.8,-4.3) -- node[rotulo, above] {\code{alerta} \code{ativo}} (10.2,-4.3);
  \draw[msg=cinza] (10.2,-4.9) -- node[rotulo, above] {``Queda no quarto 101''} (13.6,-4.9);
  % cancelamento
  \node[nota] at (0.1,-5.8) {dois toques no botão\\(até 15\,s depois)};
  \draw[msg=corvest] (0,-6.5) -- node[rotulo, above] {\code{FALL\_CANCEL}} (3.4,-6.5);
  \draw[msg=corhub] (3.4,-7.1) -- node[rotulo, above] {\code{alerta} \code{cancelado}} (6.8,-7.1);
  \draw[msg=corcentral] (6.8,-7.6) -- (10.2,-7.6);
  \draw[msg=cinza] (10.2,-8.1) -- node[rotulo, above] {``cancelada''} (13.6,-8.1);
  % metas
  \draw[decorate, decoration={brace, amplitude=4pt, mirror}, corhub] (-0.35,-1.5) -- (-0.35,-3.4)
       node[midway, left=5pt, align=right] {pico $\to$ MQTT\\$\le$ 500\,ms\\(RNF02)};
  \draw[decorate, decoration={brace, amplitude=4pt}, cinza] (14.0,-3.3) -- (14.0,-5.0)
       node[midway, right=5pt, align=left] {alerta $\to$\\Telegram\\$\le$ 5\,s (RNF04)};
\end{tikzpicture}
}
\caption{Sequência de uma queda e do cancelamento, da pulseira ao Telegram.}
\label{fig:sequencia}
\end{figure}

\subsection{MQTT: do hub à central}\label{sec:mqtt}

O hub publica tudo num broker Mosquitto~\cite{mqtt} na rede local (Tabela~\ref{tab:broker}).
Os tópicos seguem a hierarquia \code{idoso/<site>/<quarto>/<assunto>}. Assim, um
consumidor escolhe o que quer com os curingas do MQTT:
\begin{itemize}
  \item \code{idoso/butanta/quarto101/porta/banheiro}: a porta do banheiro de um quarto;
  \item \code{idoso/butanta/quarto101/\#}: tudo de um quarto;
  \item \code{idoso/+/+/alerta}: os alertas de todos os quartos, como faz o notificador.
\end{itemize}

\begin{table}[H]
\centering
\caption{O broker.}
\label{tab:broker}
\small
\begin{tabular}{@{}>{\raggedright\arraybackslash}p{3.0cm}>{\raggedright\arraybackslash}p{10.8cm}@{}}
\toprule
Item & Valor \\
\midrule
Implementação & Eclipse Mosquitto 2.0, MQTT 3.1.1 \\
Portas & 1883 (TCP: hubs e serviços), 9001 (WebSocket: navegador); 8883 com TLS quando houver nuvem \\
Autenticação & usuário e senha por cliente, com ACL por tópico: \code{hub}, \code{pulseira}, \code{telegraf}, \code{notificador}, \code{painel} \\
Keepalive & 30\,s: o \textit{Last Will} sai até $\sim$45\,s depois de a placa cair \\
\bottomrule
\end{tabular}
\end{table}

\begin{table}[H]
\centering
\caption{Tópicos publicados pelo hub (dentro de \code{idoso/<site>/<quarto>/}).}
\label{tab:topicos}
\small
\begin{tabular}{@{}lcc>{\raggedright\arraybackslash}p{4.9cm}>{\raggedright\arraybackslash}p{3.6cm}@{}}
\toprule
Tópico & QoS & Retido & Exemplo de conteúdo & Quando \\
\midrule
\code{status} & 1 & sim & \code{\{"estado":"online", "hora":true\}} & ao conectar; \textit{Last Will}: \code{offline} \\
\code{porta/<id>} & 1 & sim & \code{\{"estado":"aberta"\}} & a cada mudança \\
\code{presenca} & 1 & sim & \code{\{"estado":"detectada"\}} & a cada mudança \\
\code{temperatura} & 0 & sim & \code{\{"valor":22.5\}} & ao mudar 0,2\,°C, ou a cada 5\,min \\
\code{pulseira/<id>} & 1 & sim & \code{\{"estado":"online", "atividade":0.12\}} & no primeiro heartbeat, depois a cada 5\,min; \code{offline} após 90\,s \\
\code{alerta} & 1 & não & ver abaixo & início e fim de cada episódio \\
\code{config/aplicada} & 1 & sim & a configuração em uso & no boot e a cada \code{config/desejada} \\
\bottomrule
\end{tabular}
\end{table}

Todo JSON leva a versão do formato (\code{"v":1}) e a hora UTC da placa em milissegundos
(\code{ts}), omitida enquanto a placa não tem hora. Os estados são retidos: quem assina
depois (o painel que acabou de abrir) recebe o último valor na hora. O \code{alerta} não é
retido, porque cada mensagem é um evento; quem precisa do estado dos alertas o reconstrói
pelos episódios.

\paragraph{Envio por mudança.} Cada mensagem vira um ponto no banco da central, então o hub
não publica leitura repetida. Portas, presença e alertas já saem só quando mudam. A temperatura
e a pulseira chegam ao hub em intervalos fixos, e aí medir e enviar se separam: as regras
recebem toda leitura, e o MQTT e o CAN, só o que traz novidade.
\begin{itemize}
  \item \textbf{Temperatura}: medida a cada 60\,s e publicada quando muda 0,2\,°C ou mais em
        relação ao \emph{último valor enviado}, ou a cada 5\,min como sinal de vida. Comparar com o
        último enviado, e não com a leitura anterior, faz uma deriva lenta acabar aparecendo.
  \item \textbf{Pulseira}: o heartbeat continua a cada 30\,s pelo ESP-NOW, porque é ele que dá o
        \code{offline} em 90\,s; mas \code{pulseira/<id>} sai no primeiro heartbeat e depois como
        um resumo a cada 5\,min, com a atividade calculada no período. O \code{offline} sai na hora.
\end{itemize}
Com o quarto estável, cada um passa de 2.880 para cerca de 300 pontos por dia; uma mudança
real de temperatura chega ao painel na hora. A regra fica num componente comum, em C puro e
testado no computador (\code{report}).

\noindent Um alerta de queda:
\begin{lstlisting}[language=json]
{"v":1, "id":"a1b2c3-42", "tipo":"queda", "estado":"ativo",
 "origem":"pulseira/a1b2c3", "detalhe":{"a2":7.88, "prob":0.97}, "ts":1792497600000}
\end{lstlisting}

O \code{id} identifica o episódio e se repete no \code{resolvido} ou \code{cancelado}
correspondente. Com QoS 1, uma mensagem pode chegar duas vezes; quem consome descarta
\code{id} e \code{estado} repetidos. O \code{detalhe} da queda traz o $\amag^2$ no pico e a
probabilidade que o classificador da pulseira deu. O pedido de ajuda é um alerta do tipo
\code{ajuda}, com o \code{id} formado do mesmo jeito e o \code{detalhe} vazio.

\paragraph{Reserva pela pulseira.} Com a opção MQTT ligada, a pulseira se associa ao mesmo
Wi-Fi e, se o hub não confirmar uma queda ou um pedido de ajuda, publica o alerta ela mesma,
com o mesmo \code{id};
a central recebe uma notificação só. Custa bateria, por isso vem desligado.

\subsection{Segurança}

\begin{itemize}
  \item ESP-NOW cifrado (PMK e LMK), e o hub aceita só a pulseira pareada: um ESP32 qualquer
        por perto não consegue mandar uma queda falsa.
  \item MQTT com usuário e senha por cliente e ACL por tópico: o notificador só lê alertas, a
        pulseira só publica a própria telemetria e o alerta, e só o painel escreve a
        configuração.
  \item Senhas, chaves e o \code{sdkconfig} ficam fora do repositório.
  \item Na rede local, o MQTT ainda é sem TLS; o TLS entra quando houver ponte para a nuvem.
\end{itemize}

\begin{figure}[H]
\centering
\fotoouespaco{serial_espnow}{0.85\linewidth}{monitor serial do hub recebendo o heartbeat e uma queda da pulseira por ESP-NOW}
\caption{Hub recebendo a pulseira por ESP-NOW.}
\label{fig:serial_espnow}
\end{figure}
